\chapter{Lesson 1: The Weekend}

The first new lesson of the term is about what we did last weekend, which means the
\textbf{past tense} is the main grammar point. Monday classes for the rest of the term
opened with the question \ko{주말 잘 보냈어요?} Jumal jal bonaesseoyo? (Did you have a
good weekend?), so it pays to have a few sentences ready.

\section{Vocabulary}

\begin{table}[H]
\centering
\begin{tabular}{lll}
\toprule
Korean & RR & English \\
\midrule
\ko{주말} & jumal & weekend \\
\ko{평일} & pyeongil & weekday \\
\ko{지난 주말} & jinan jumal & last weekend \\
\ko{청소하다} & cheongsohada & to clean \\
\ko{빨래하다} & ppallaehada & to do laundry \\
\ko{요리하다} & yorihada & to cook \\
\ko{운동하다} & undonghada & to exercise \\
\ko{등산하다} & deungsanhada & to go hiking (up a mountain) \\
\ko{산책하다} & sanchaekhada & to take a walk \\
\ko{쇼핑하다} & syopinghada & to go shopping \\
\ko{아르바이트하다} & areubaiteuhada & to work part-time \\
\ko{친구를 만나다} & chingureul mannada & to meet a friend \\
\ko{영화를 보다} & yeonghwareul boda & to watch a movie \\
\ko{늦잠을 자다} & neutjameul jada & to sleep in \\
\ko{쉬다} & swida & to rest \\
\ko{놀다} & nolda & to hang out, to play \\
\ko{보내다} & bonaeda & to spend (time); to send \\
\ko{재미있다} & jaemiitda & to be fun, interesting \\
\ko{피곤하다} & pigonhada & to be tired \\
\bottomrule
\end{tabular}
\end{table}

\subsection{Days of the week and time words}

\begin{table}[H]
\centering
\begin{tabular}{llll}
\toprule
Korean & RR & English & Element \\
\midrule
\ko{월요일} & woryoil & Monday & moon (\ko{월}) \\
\ko{화요일} & hwayoil & Tuesday & fire (\ko{화}) \\
\ko{수요일} & suyoil & Wednesday & water (\ko{수}) \\
\ko{목요일} & mogyoil & Thursday & wood (\ko{목}) \\
\ko{금요일} & geumyoil & Friday & metal, gold (\ko{금}) \\
\ko{토요일} & toyoil & Saturday & earth (\ko{토}) \\
\ko{일요일} & iryoil & Sunday & sun (\ko{일}) \\
\bottomrule
\end{tabular}
\end{table}

\begin{center}
\begin{tikzpicture}[x=2.3cm, every node/.style={font=\small}]
  \draw[-{Stealth}, thick] (-0.3,0) -- (4.4,0) node[right] {time};
  \foreach \x/\kw/\rr/\ee in {0/그저께/geujeokke/2 days ago, 1/어제/eoje/yesterday, 2/오늘/oneul/today, 3/내일/naeil/tomorrow, 4/모레/more/in 2 days} {
    \draw (\x,0.1) -- (\x,-0.1);
    \node[above] at (\x,0.15) {\ko{\kw}};
    \node[below, font=\scriptsize, align=center] at (\x,-0.15) {\rr\\ \ee};
  }
  \draw[-{Stealth}, thick] (-0.3,-2.2) -- (4.4,-2.2) node[right] {weeks};
  \foreach \x/\kw/\rr in {1/지난주/jinanju, 2/이번 주/ibeon ju, 3/다음 주/daeum ju} {
    \draw (\x,-2.1) -- (\x,-2.3);
    \node[above] at (\x,-2.05) {\ko{\kw}};
    \node[below, font=\scriptsize] at (\x,-2.35) {\rr};
  }
  \fill[red!60] (2,0) circle (2.5pt);
  \fill[red!60] (2,-2.2) circle (2.5pt);
  \node[font=\scriptsize, red!60!black] at (2,0.85) {now (\ko{지금})};
\end{tikzpicture}
\end{center}

Days of the week and clock times take the time particle \ko{에}: \ko{토요일에} toyoire
(on Saturday), \ko{세 시에} se sie (at three). Relative time words do \emph{not}:
\ko{어제}, \ko{오늘}, \ko{내일}, \ko{지금}, \ko{지난주} are used bare. (\ko{어제에} is a
classic beginner mistake.) \ko{주말에} jumare (on the weekend) does take \ko{에}.

\section{Grammar 1: Past tense \ko{\til 았/었어요}}

The past tense is built by inserting \ko{았} or \ko{었} between the stem and the
polite ending. The choice follows exactly the same vowel-harmony rule as the present
tense, so if you can form the present you can form the past: take the present
\ko{해요} form, drop \ko{요}, and add \ko{ㅆ어요}.

\begin{center}
\begin{tikzpicture}[every node/.style={font=\small},
    b/.style={draw, rounded corners, minimum height=0.85cm, minimum width=1.4cm}]
  \node[b, fill=gray!10] (s) at (0,0) {\ko{먹}};
  \node[b, fill=orange!15] (t) at (2.0,0) {\ko{었}};
  \node[b, fill=blue!10] (e) at (4.0,0) {\ko{어요}};
  \node at (1.0,0) {+}; \node at (3.0,0) {+};
  \node at (5.3,0) {$=$};
  \node[b, fill=green!10, minimum width=2cm] at (6.7,0) {\ko{먹었어요}};
  \node[font=\scriptsize, below=0.2cm of s] {stem};
  \node[font=\scriptsize, below=0.2cm of t] {past marker};
  \node[font=\scriptsize, below=0.2cm of e] {polite ending};
  \node[font=\scriptsize] at (6.7,-0.75) {meogeosseoyo, ``ate''};
\end{tikzpicture}
\end{center}

\begin{table}[H]
\centering
\begin{tabular}{lllll}
\toprule
Dictionary & Present & Past & RR & Meaning \\
\midrule
\ko{가다} & \ko{가요} & \ko{갔어요} & gasseoyo & went \\
\ko{오다} & \ko{와요} & \ko{왔어요} & wasseoyo & came \\
\ko{보다} & \ko{봐요} & \ko{봤어요} & bwasseoyo & saw, watched \\
\ko{만나다} & \ko{만나요} & \ko{만났어요} & mannasseoyo & met \\
\ko{놀다} & \ko{놀아요} & \ko{놀았어요} & norasseoyo & hung out \\
\ko{먹다} & \ko{먹어요} & \ko{먹었어요} & meogeosseoyo & ate \\
\ko{쉬다} & \ko{쉬어요} & \ko{쉬었어요} & swieosseoyo & rested \\
\ko{마시다} & \ko{마셔요} & \ko{마셨어요} & masyeosseoyo & drank \\
\ko{배우다} & \ko{배워요} & \ko{배웠어요} & baewosseoyo & learned \\
\ko{보내다} & \ko{보내요} & \ko{보냈어요} & bonaesseoyo & spent (time) \\
\ko{쓰다} & \ko{써요} & \ko{썼어요} & sseosseoyo & wrote \\
\ko{하다} & \ko{해요} & \ko{했어요} & haesseoyo & did \\
\ko{재미있다} & \ko{재미있어요} & \ko{재미있었어요} & jaemiisseosseoyo & was fun \\
\bottomrule
\end{tabular}
\end{table}

\begin{important}[Past of the copula]
\ko{이다} becomes \ko{이었어요} after a consonant and \ko{였어요} after a vowel:
\ko{어제는 제 생일이었어요.} Eojeneun je saengirieosseoyo. (Yesterday was my birthday.)
\ko{그 사람은 가수였어요.} Geu sarameun gasuyeosseoyo. (That person was a singer.) The
negative is \ko{아니었어요}: \ko{학생이 아니었어요} haksaengi anieosseoyo (wasn't a
student).
\end{important}

\begin{example}[Talking about last weekend]
Some sentences I prepared for the Monday warm-up:
\begin{itemize}
  \item \ko{토요일에 늦잠을 잤어요.} Toyoire neutjameul jasseoyo. (I slept in on
        Saturday.)
  \item \ko{오후에 도서관에서 숙제를 했어요.} Ohue doseogwaneseo sukjereul haesseoyo. (In
        the afternoon I did homework at the library.)
  \item \ko{일요일에는 친구를 만났어요. 같이 영화를 봤어요.} Iryoireneun chingureul
        mannasseoyo. Gachi yeonghwareul bwasseoyo. (On Sunday I met a friend. We watched
        a movie together.)
  \item \ko{영화가 아주 재미있었어요.} Yeonghwaga aju jaemiisseosseoyo. (The movie was
        really fun.)
\end{itemize}
\end{example}

\subsection{Negative past}

There are two ways to negate, and both carry over to the past:
\begin{itemize}
  \item \textbf{Short negation} \ko{안} + verb: \ko{안 갔어요} an gasseoyo (didn't go).
        With \ko{noun + 하다} verbs, \ko{안} goes right before \ko{하다}:
        \ko{운동 안 했어요} undong an haesseoyo, \emph{not} \ko{안 운동했어요}.
  \item \textbf{Long negation} stem + \ko{\til 지 않았어요}: \ko{가지 않았어요} gaji
        anasseoyo. Slightly more formal; common in writing.
\end{itemize}

\begin{example}[Negative answers]
Answering a question with the opposite of what was asked:
\begin{itemize}
  \item \ko{주말에 청소했어요? --- 아니요, 청소 안 했어요. 너무 피곤했어요.} Jumare
        cheongsohaesseoyo? --- Aniyo, cheongso an haesseoyo. Neomu pigonhaesseoyo. (Did
        you clean on the weekend? --- No, I didn't. I was too tired.)
  \item \ko{어제 아침을 먹지 않았어요.} Eoje achimeul meokji anasseoyo. (I didn't eat
        breakfast yesterday.)
\end{itemize}
\end{example}

\section{Grammar 2: ``and'' and ``with'' for nouns}

Korean has several particles meaning ``and'' (joining nouns) or ``with'' (a companion).
They differ in register, not meaning.

\begin{table}[H]
\centering
\begin{tabular}{llll}
\toprule
Particle & After consonant / vowel & Register & Example \\
\midrule
\ko{하고} & \ko{하고} / \ko{하고} & spoken, neutral & \ko{빵하고 우유} \\
\ko{와/과} & \ko{과} / \ko{와} & written, formal & \ko{책과 공책}, \ko{커피와 차} \\
\ko{(이)랑} & \ko{이랑} / \ko{랑} & spoken, casual & \ko{동생이랑}, \ko{친구랑} \\
\bottomrule
\end{tabular}
\end{table}

Notice that \ko{와/과} runs ``backwards'' compared with most particles: \ko{과} after a
consonant, \ko{와} after a vowel. To say ``together with'', add \ko{같이} gachi or
\ko{함께} hamkke: \ko{친구하고 같이 등산했어요.} Chinguhago gachi deungsanhaesseoyo. (I
went hiking with a friend.)

\begin{remark}
These particles join \emph{nouns} only. To join two \emph{clauses} (``I cleaned and then
rested'') use \ko{\til 고} (Lesson 3) or the sentence-initial connector \ko{그리고}.
\end{remark}

\section{Sentence connectors}

Paragraph-length answers sound much better with connecting adverbs at the start of
sentences.

\begin{table}[H]
\centering
\begin{tabular}{lll}
\toprule
Korean & RR & Meaning \\
\midrule
\ko{그리고} & geurigo & and, and then \\
\ko{그래서} & geuraeseo & so, therefore \\
\ko{그런데} & geureonde & but, by the way (topic shift) \\
\ko{그렇지만} & geureochiman & however \\
\ko{그 다음에} & geu daeume & after that \\
\bottomrule
\end{tabular}
\end{table}

\begin{example}[A short diary entry]
A short paragraph in the style of the weekly journal assignment:
\begin{quote}
\ko{지난 주말에 날씨가 좋았어요. 그래서 친구랑 공원에서 산책했어요. 그 다음에 식당에서
비빔밥을 먹었어요. 그런데 저녁에 비가 왔어요. 그래서 집에서 쉬었어요.}
\end{quote}
Jinan jumare nalssiga joasseoyo. Geuraeseo chingurang gongwoneseo sanchaekhaesseoyo. Geu
daeume sikdangeseo bibimbabeul meogeosseoyo. Geureonde jeonyeoge biga wasseoyo.
Geuraeseo jibeseo swieosseoyo. (Last weekend the weather was nice, so I took a walk in
the park with a friend. After that I ate bibimbap at a restaurant. But in the evening it
rained, so I rested at home.)
\end{example}

\section{Pronunciation notes}

\begin{itemize}
  \item The past marker \ko{ㅆ} links onto \ko{어요}: \ko{갔어요} is [\ko{가써요}]. When
        followed by a consonant it neutralizes to [t] and tensifies the next consonant:
        \ko{갔지만} [\ko{갇찌만}] gatjiman.
  \item \ko{했어요} is often reduced in fast speech to something like [\ko{해써요}] with
        a very short first vowel.
  \item \ko{같이} is [\ko{가치}] by palatalization (Chapter 1).
\end{itemize}

\begin{aside}
\textbf{Weekend culture.} Our instructor described typical university-student weekends in
Seoul: \ko{노래방} noraebang (karaoke rooms rented by the hour), \ko{PC방} PC-bang
(internet caf\'es for gaming), and \ko{찜질방} jjimjilbang (24-hour bathhouses with
heated rooms, where families and friends hang out and even sleep over). Hiking
(\ko{등산}) is hugely popular with middle-aged Koreans, who are known for wearing full
professional hiking gear even on easy city trails.
\end{aside}

\section*{Review summary}
\begin{itemize}
  \item Past tense: stem + \ko{았} (last vowel \ko{ㅏ/ㅗ}) or \ko{었} (otherwise) +
        \ko{어요}; \ko{하다} $\to$ \ko{했어요}; contractions as in the present.
  \item Copula past: \ko{이었어요} / \ko{였어요}; negative \ko{아니었어요}.
  \item Negation: \ko{안} + verb (before \ko{하다} in noun-\ko{하다} verbs) or
        \ko{\til 지 않았어요}.
  \item ``And/with'' for nouns: \ko{하고} (neutral), \ko{와/과} (written),
        \ko{(이)랑} (casual); add \ko{같이} for ``together''.
  \item Days of the week and clock times take \ko{에}; \ko{어제, 오늘, 내일} do not.
\end{itemize}
